\documentclass[11pt,a4paper]{article}
\usepackage{amsmath,amssymb,amsthm}
\usepackage{alphalph}
\usepackage[margin=2.4cm]{geometry}
\emergencystretch=2em
\newtheorem{theorem}{Theorem}[section]
\newtheorem{lemma}[theorem]{Lemma}
\newtheorem{proposition}[theorem]{Proposition}
\newtheorem{corollary}[theorem]{Corollary}
\theoremstyle{definition}
\newtheorem{definition}[theorem]{Definition}
\newtheorem{convention}[theorem]{Convention}
\newtheorem{remark}[theorem]{Remark}
\providecommand{\fint}{\ensuremath{\mathchoice{\kern.18em\overline{\kern-.18em\int}}{\overline{\int}}{\overline{\int}}{\overline{\int}}}}
\providecommand{\ddc}{dd^c}
\providecommand{\href}[2]{#2}
\newcommand{\lean}[1]{}
\newcommand{\uses}[1]{}
\newcommand{\source}[1]{}
\newcommand{\dcref}[1]{\textup{(#1)}}
\newcommand{\leanok}{}
\newcommand{\mathlibok}{}
\newcommand{\notready}{}
\title{Hironaka: dimension components, transversality and divisor multiplicities}
\author{}
\date{11 October 2026}
\begin{document}
\maketitle
\section{Prerequisites from the original API}
These results use the embedded resolution theorems already proved in the library.
The following pointers identify the earlier results referred to in the proofs below.
\subsection{Pure-dimensional embedded resolution}\label{hir:embedded-resolution-pure-dimensional}
See \texttt{Hironaka.exists\_embeddedResolution\_pure\_of\_isCompact} and its compact-ambient and relatively compact variants in \texttt{API} and \texttt{API2}.
\subsection{Embedded resolution with an SNC boundary}\label{hir:embedded-resolution-boundary}
See \texttt{Hironaka.exists\_embeddedResolution\_boundary\_of\_isCompact} and the corresponding variants in \texttt{API2}.
\subsection{Embedded resolution in mixed dimensions}\label{hir:embedded-resolution-analytic}
See \texttt{Hironaka.exists\_embeddedResolution\_analytic\_of\_isCompact} and its boundary form in \texttt{API3}.
\section{Dimension parts and genuine global irreducible components}

Let $A$ be a closed analytic subset of a complex manifold $M$ of dimension $n$.
The global irreducible components are the closures of the connected components
of the regular locus, as in \cite{demailly-cadg}, II.5.3. A set has pure
geometric dimension $k$ when every one of its regular points has dimension $k$.
This characterization uses actual local manifold dimensions; it assigns no
extraneous dimension certificate. Write $A^{(k)}=\overline{\operatorname{Reg}_k(A)}$.

\begin{lemma}[Pure geometric dimension of a global component]
\label{hir:global-component-pure-geometric-dimension}
\uses{ana:regular-component,ana:pure-dimension,rem:global-components}
\source{demailly-cadg}
\dcref{II.5.3}
\lean{Hironaka.HironakaDimensionGlobalComponents.isPureDim_closure_component,Hironaka.HironakaDimensionGlobalComponents.exists_unique_isPureDim_of_mem_globalComponents}
\leanok
Every global irreducible component of $A$ has a unique pure geometric dimension
$k\leq n$.
\end{lemma}
\begin{proof}
A global component $B$ is the closure of the connected component $C$ through
some regular point $c$ of $A$, say of dimension $k$. All points of $C$ are
regular points of $A$ of dimension $k$, and they are also regular points of
$B$ of dimension $k$. In particular $B$ has a regular point of dimension $k$.
The closure $B$ is an irreducible analytic set. Pure dimensionality of an
irreducible analytic set therefore implies that every regular point of $B$
has dimension $k$. The local dimension at $c$ is at most $n$. If $B$ also has
pure dimension $j$, its regular point $c$ gives $j=k$, proving uniqueness.
\end{proof}

\begin{lemma}[Closed unions of global components]
\label{hir:closed-global-component-subunion}
\uses{rem:global-components}
\source{demailly-cadg}
\dcref{II.5.3}
\lean{Hironaka.HironakaDimensionGlobalComponents.isClosed_sUnion_globalComponents_subfamily}
\leanok
The union of any subfamily of the global irreducible components of $A$ is
closed in $M$.
\end{lemma}
\begin{proof}
Every global component is closed, and their family is locally finite. Any
subfamily remains locally finite. A locally finite union of closed sets is
closed: at a point outside the union, choose a neighborhood meeting only
finitely many members, and remove those finitely many closed sets from the
neighborhood. The resulting neighborhood misses the entire union.
\end{proof}

\begin{theorem}[The irreducible-component description of dimension parts]
\label{hir:dimension-part-global-component-union}
\uses{hir:global-component-pure-geometric-dimension,hir:closed-global-component-subunion,hir:embedded-resolution-dimension-parts}
\source{demailly-cadg}
\dcref{II.5.3}
\lean{Hironaka.HironakaDimensionGlobalComponents.dimensionPart_eq_sUnion_globalComponents}
\leanok
For every $k\in\mathbb N$, the dimension part $A^{(k)}$ is exactly the union
of the global irreducible components of $A$ of pure geometric dimension $k$.
No compactness, empty-interior, or positive-dimension assumption is required.
\end{theorem}
\begin{proof}
Let $S$ be that union. If $c$ is a regular point of $A$ of dimension $k$,
the closure of its connected component in the regular locus is a global
component of pure dimension $k$ containing $c$. Hence
$\operatorname{Reg}_k(A)\subset S$. The preceding lemma makes $S$ closed,
so $A^{(k)}\subset S$.

Conversely, let $B$ be a global component of pure dimension $k$, and write
$B=\overline C$ for the regular-locus component through a regular point $c$
of $A$ of dimension $j$. That same point is regular of dimension $j$ on $B$,
so pure dimension $k$ gives $j=k$. Every point of $C$ is then regular of
dimension $k$ on $A$, and therefore
$B=\overline C\subset\overline{\operatorname{Reg}_k(A)}=A^{(k)}$.
Taking the union gives $S\subset A^{(k)}$ and the claimed equality.
\end{proof}

\begin{lemma}[Geometric purity of the dimension part]
\label{hir:dimension-part-geometric-purity}
\uses{hir:embedded-resolution-dimension-parts,lag:regular-point-basics}
\lean{Hironaka.HironakaDimensionGlobalComponents.isPureDim_dimensionPart}
\leanok
Every regular point of $A^{(k)}$ has local manifold dimension $k$.
\end{lemma}
\begin{proof}
The regular points of dimension $k$ on $A^{(k)}$ are dense in $A^{(k)}$.
At a regular point $y$ of dimension $j$ on that set, choose a neighborhood
where all points of $A^{(k)}$ are regular of dimension $j$. Density supplies
a regular point of dimension $k$ inside this neighborhood. Uniqueness of
local manifold dimension at that point gives $j=k$.
\end{proof}


\begin{lemma}[Component membership in a dimension part]
\label{hir:global-component-dimension-part-criterion}
\uses{hir:global-component-pure-geometric-dimension,hir:dimension-part-global-component-union,lag:regular-point-basics}
\lean{Hironaka.HironakaDimensionGlobalComponents.isPureDim_iff_subset_dimensionPart_of_mem_globalComponents}
\leanok
A global irreducible component $B$ of $A$ is contained in $A^{(k)}$ if and
only if it has pure geometric dimension $k$.
\end{lemma}
\begin{proof}
The reverse implication follows from the union description. For the forward
implication, write $B=\overline C$ where $C$ is the regular-locus component
through a regular point $c$ of $A$ of dimension $j$. Since $c\in B\subset
A^{(k)}$, every neighborhood of $c$ meets $\operatorname{Reg}_k(A)$. Take a
neighborhood where all points of $A$ are regular of dimension $j$. A point
of $\operatorname{Reg}_k(A)$ in it has both dimensions $k$ and $j$, so $j=k$.
The component-purity lemma gives pure dimension $k$ for $B$.
\end{proof}

\begin{theorem}[Pure local decompositions of a dimension part]
\label{hir:dimension-part-local-decomposition}
\uses{ana:dimension-strata,lag:local-statement,hir:embedded-resolution-dimension-parts}
\source{demailly-cadg}
\dcref{II.4.24, II.5.3}
\lean{Hironaka.HironakaDimensionGlobalComponents.localDecomposition_dimensionPart,Hironaka.HironakaDimensionGlobalComponents.exists_localDecomposition_dimensionPart}
\leanok
At every point of $A^{(k)}$ there is a finite local analytic decomposition
of $A^{(k)}$ whose pieces all have dimension $k$. More precisely, from any
local analytic decomposition $A\cap U=\bigcup_i A_i$ with piece dimensions
$d_i$, retain exactly the indices with $d_i=k$. They give such a decomposition
on the same $U$, with the same holomorphic equations for each retained piece,
dense regular loci and pairwise relatively thin intersections.
\end{theorem}
\begin{proof}
The local stratum-closure theorem identifies
$A^{(k)}\cap U=\bigcup_{d_i=k} A_i$. The selected index set is finite, so
choose a bijection from a finite consecutive index set onto it. Transport
the retained pieces and their equations along this bijection. Their union
is the displayed dimension part. The neighborhood is unchanged, hence
remains open, contains the chosen point and lies in its preferred chart.
Holomorphicity of every equation, the equation descriptions of each piece,
purity and density of its regular locus remain exactly the corresponding
properties of the original piece. The assigned dimension is $k$ because
the index is selected. Distinct selected indices remain distinct original
indices by injectivity of the bijection, so the relatively thin-intersection
property is inherited as well. Finally any point of $A^{(k)}$ lies in $A$;
existence of the original analytic local decomposition at that point supplies
the asserted decomposition of $A^{(k)}$.
\end{proof}

\section{Transversality of the resolved set and the accumulated divisor}

\begin{lemma}[Density excludes a containing coordinate hyperplane]
\label{hir:nc-density-transverse}
\uses{hir:snc}
\lean{Hironaka.HironakaNormalCrossingsTransversality.isTransverse_of_hasNormalCrossings_of_subset_closure_diff}
\leanok
Let \(C\) be a subset of a complex manifold and \(E=[H_1,\ldots,H_s]\) a finite family of subsets. Assume that \(C\) has only normal crossings with \(E\) and
\[
 C\subset\overline{C\setminus|E|},\qquad |E|=\bigcup_i H_i.
\]
Then \(E\) is transverse to \(C\): in the normal crossings coordinates at every point of \(C\), the coordinate defining any member of \(E\) through that point is not among the coordinates defining \(C\).
\end{lemma}
\begin{proof}
Fix \(x\in C\) and normal crossings coordinates on an open neighbourhood \(U\) of \(x\). Write \(C\cap U=\{z_j=0:j\in S_0\}\) and, for every member \(H_i\) through \(x\), write \(H_i\cap U=\{z_{a_i}=0\}\). If \(a_i\in S_0\), then \(C\cap U\subset H_i\). Density gives a point of \((C\setminus|E|)\cap U\), which would belong to \(H_i\subset|E|\), a contradiction. Thus \(a_i\notin S_0\) for every such member, which is precisely transversality. No connectedness or irreducibility is required.
\end{proof}

\begin{lemma}[The initial boundary is the only additional obstruction]
\label{hir:strict-transform-initial-boundary-transverse}
\uses{hir:nc-density-transverse,ana:strict-transform-composition,hir:admissible-calculus}
\lean{Hironaka.HironakaNormalCrossingsTransversality.isTransverse_strictTransform_of_support_subset,Hironaka.HironakaNormalCrossingsTransversality.isTransverse_strictTransform_of_admissible_of_preimage_subset_exceptional}
\leanok
Let \(f:Y\to X\) be holomorphic, \(A\) closed, \(B\subset X\), and \(E\) a finite family on \(Y\) such that
\[
 A\subset\overline{A\setminus B},\qquad
 |E|\subset\operatorname{Exc}(f)\cup f^{-1}(B).
\]
If the strict transform \(A'\) has only normal crossings with \(E\), then \(E\) is transverse to \(A'\). In particular, this applies to an admissible sequence of \((X,E_0)\) with centres over \(G\cup(A\cap|E_0|)\), whenever \(f^{-1}(G)\subset\operatorname{Exc}(f)\) and \(A\subset\overline{A\setminus|E_0|}\).
\end{lemma}
\begin{proof}
Put \(D=f^{-1}(A)\setminus\operatorname{Exc}(f)\), so \(A'=\overline D\). Fix \(y\in D\) and a neighbourhood \(V\) of \(y\). The exceptional locus has open complement, and \(f\) is a local biholomorphism at \(y\). Consequently \(f(V\setminus\operatorname{Exc}(f))\) is a neighbourhood of \(f(y)\). Density off \(B\) provides a point \(f(w)\in A\setminus B\) with \(w\in V\setminus\operatorname{Exc}(f)\). Then \(w\in D\subset A'\), and the support inclusion gives \(w\notin|E|\). Every neighbourhood of \(y\) therefore meets \(A'\setminus|E|\). Taking closures gives \(A'\subset\overline{A'\setminus|E|}\), and Lemma~\ref{hir:nc-density-transverse} applies.

For the final assertion, the accumulated support lies in \(f^{-1}(|E_0|\cup G)\), since \(A\cap|E_0|\subset|E_0|\). The assumed inclusion over \(G\) puts that support inside \(f^{-1}(|E_0|)\cup\operatorname{Exc}(f)\), giving the required hypotheses.
\end{proof}

\begin{lemma}[Transversality after an admissible sequence]
\label{hir:strict-transform-density-transverse}
\uses{hir:nc-density-transverse,hir:admissible-calculus}
\lean{Hironaka.HironakaNormalCrossingsTransversality.isTransverse_strictTransform_of_admissible}
\leanok
Let \(f:Y\to X\) be an admissible blowup sequence of \((X,E_0)\), with centres over \(F\) and accumulated divisor \(E\). Suppose the members of \(E_0\) are closed, \(A\subset X\) is closed,
\[
 A\subset\overline{A\setminus(|E_0|\cup F)},
\]
and the strict transform \(A'\) has only normal crossings with \(E\). Then \(E\) is transverse to \(A'\).
\end{lemma}
\begin{proof}
The support of the accumulated divisor lies in \(f^{-1}(|E_0|\cup F)\). Density survives passage to the strict transform, so
\[
 A'\subset\overline{A'\setminus f^{-1}(|E_0|\cup F)}
 \subset\overline{A'\setminus|E|}.
\]
Apply Lemma~\ref{hir:nc-density-transverse}.
\end{proof}

\begin{theorem}[Embedded resolution with a transverse boundary]
\label{hir:embedded-resolution-transverse-boundary}
\uses{hir:strict-transform-density-transverse,hir:embedded-resolution-pure-dimensional,hir:embedded-resolution-boundary}
\source{bierstone-milman97-invent-math:page-0009}
\lean{Hironaka.HironakaNormalCrossingsTransversality.exists_embeddedResolution_pure_transverse_of_isCompact,Hironaka.HironakaNormalCrossingsTransversality.exists_embeddedResolution_boundary_transverse_of_isCompact,Hironaka.HironakaNormalCrossingsTransversality.exists_embeddedResolution_pure_transverse,Hironaka.HironakaNormalCrossingsTransversality.exists_embeddedResolution_boundary_transverse}
\leanok
Let \(W\) be a compact analytic subset of a complex manifold of pure dimension \(d<n\). Its embedded resolution can be chosen with all the conclusions of Theorem~\ref{hir:embedded-resolution-pure-dimensional}, and with the accumulated divisor transverse to the smooth strict transform. If \(E_0\) is a simple normal crossings divisor and \(W\subset\overline{W\setminus|E_0|}\), the same strengthening holds for Theorem~\ref{hir:embedded-resolution-boundary}, retaining its centres over \(W_{\mathrm{sing}}\cup(W\cap|E_0|)\). These statements apply in particular to every such \(W\) in a compact ambient manifold.
\end{theorem}
\begin{proof}
For the empty initial divisor the regular points of dimension \(d\) are dense in \(W\) and avoid \(W_{\mathrm{sing}}=W\setminus\operatorname{Reg}_d(W)\). Therefore \(W\subset\overline{W\setminus W_{\mathrm{sing}}}\), and Lemma~\ref{hir:strict-transform-density-transverse} applies to the sequence of Theorem~\ref{hir:embedded-resolution-pure-dimensional}.

For the boundary form, put \(B=|E_0|\), which is closed. We claim that \(W\subset\overline{W\setminus(W_{\mathrm{sing}}\cup B)}\). Given \(x\in W\) and an open neighbourhood \(U\) of \(x\), density off \(B\) gives \(y\in(W\setminus B)\cap U\). The open set \(U\setminus B\) contains \(y\), so density of \(\operatorname{Reg}_d(W)\) gives \(z\in\operatorname{Reg}_d(W)\cap(U\setminus B)\). This proves the claim. Since
\[
 B\cup\bigl(W_{\mathrm{sing}}\cup(W\cap B)\bigr)
 =W_{\mathrm{sing}}\cup B,
\]
Lemma~\ref{hir:strict-transform-density-transverse} applies to the existing boundary resolution. Its conclusion strengthens the normal crossings conclusion without changing the support conditions, dimensions, or position of the centres.
\end{proof}

\begin{theorem}[Transverse resolution without equidimensionality]
\label{hir:embedded-resolution-analytic-transverse}
\uses{hir:strict-transform-initial-boundary-transverse,hir:embedded-resolution-analytic}
\source{bierstone-milman97-invent-math:page-0009}
\lean{Hironaka.HironakaNormalCrossingsTransversality.exists_embeddedResolution_analytic_transverse_of_isCompact,Hironaka.HironakaNormalCrossingsTransversality.exists_embeddedResolution_analytic_boundary_transverse_of_isCompact}
\leanok
Let \(W\) be compact analytic with empty interior, allowing components of different dimensions. The embedded resolutions of Theorem~\ref{hir:embedded-resolution-analytic}, with or without an initial simple normal crossings divisor \(E_0\) satisfying \(W\subset\overline{W\setminus|E_0|}\), can be chosen with the accumulated divisor transverse to \(W'\). All dimension-part, disjointness, exceptional-support and isomorphism conclusions are retained.
\end{theorem}
\begin{proof}
The boundary theorem supplies an admissible sequence with centres over \(W_{\mathrm{sing}}\cup(W\cap|E_0|)\), with \(f^{-1}(W_{\mathrm{sing}})\subset\operatorname{Exc}(f)\), and with \(W'\) having normal crossings with \(E\). Apply Lemma~\ref{hir:strict-transform-initial-boundary-transverse} with \(G=W_{\mathrm{sing}}\). This argument does not use equidimensionality. For the empty initial divisor, the old theorem gives \(|E|=\operatorname{Exc}(f)\). Apply the same lemma with \(B=\varnothing\). These operations retain the original witnesses and every stated conclusion about their dimension parts.
\end{proof}

\begin{theorem}[Transverse embedded resolution on relatively compact open sets]
\label{hir:embedded-resolution-transverse-local}
\uses{hir:strict-transform-initial-boundary-transverse,hir:embedded-resolution-boundary,hir:embedded-resolution-pure-dimensional,hir:boundary-relatively-compact}
\lean{Hironaka.HironakaNormalCrossingsTransversality.exists_embeddedResolution_pure_transverse_of_subset_isCompact,Hironaka.HironakaNormalCrossingsTransversality.exists_embeddedResolution_boundary_transverse_of_subset_isCompact,Hironaka.HironakaNormalCrossingsTransversality.exists_embeddedResolution_boundary_transverse_local}
\leanok
The transverse conclusions of Theorem~\ref{hir:embedded-resolution-transverse-boundary} also hold on any open set \(U\subset X\) contained in a compact set, without requiring \(W\) compact. In particular, every compact subset \(K\subset X\) has an open neighbourhood \(U\) with compact closure on which the boundary resolution has this transverse conclusion.
\end{theorem}
\begin{proof}
Use the existing embedded resolutions on \(U\). In the empty-boundary case the accumulated support is exactly the exceptional locus, so Lemma~\ref{hir:strict-transform-initial-boundary-transverse} applies with \(B=\varnothing\).

For the boundary case put \(W_U=W\cap U\) and \(E_{0,U}=E_0|_U\). The inclusion \(U\hookrightarrow X\) is an open embedding, so pulling back the closure of any set equals the closure of its pullback. Hence density of \(W\setminus|E_0|\) in \(W\) gives density of \(W_U\setminus|E_{0,U}|\) in \(W_U\). The old theorem supplies a sequence with centres over \((W_{\mathrm{sing}}\cap U)\cup(W_U\cap|E_{0,U}|)\), and the entire preimage of \(W_{\mathrm{sing}}\cap U\) is exceptional. Lemma~\ref{hir:strict-transform-initial-boundary-transverse} therefore gives transversality to the accumulated divisor, retaining every existing conclusion of the local resolution.

A complex manifold is locally compact. Cover a compact \(K\) by finitely many open neighbourhoods with compact closures and take their union \(U\). Its closure is contained in the finite union of those compact closures, so it is compact. Applying the preceding result on \(U\subset\overline U\) completes the neighbourhood form.
\end{proof}

\begin{remark}
The density argument supplies the noncontainment needed for simultaneous normal crossings in Bierstone--Milman's Theorem 1.6(2). It does not change the locus over which the boundary resolution may blow up, and does not treat an initial divisor with general, nonsimple normal crossings.
\end{remark}

\section{Divisor exponents of genuine ideal powers}
\label{sec:hironaka-divisor-extension}

The monomial factor in W{\l}odarczyk's Step~2 is a product of principal
divisor ideals, with multiplicities.  The following local algebra records
how those multiplicities change when an ideal is replaced by a genuine
ideal power.  It is an extension of the order and divisor-exponent calculus,
not a construction of the global decomposition of every principalizing
divisor into the members of a prescribed simple normal crossings family.

\begin{lemma}[Successor thresholds for ideal powers]
\label{hir:divisor-power-successor-threshold}
\uses{hir:order-power,hir:order-basic,hir:cosupport}
\lean{Hironaka.HironakaDivisorPowers.cosupport_idealPow_mul_succ_eq}
\leanok
Let $\mathfrak a$ be an ideal of holomorphic germs on a complex manifold
$M$, and let $k\geq 1$ and $e\geq 0$ be integers.  Then
\[
 \operatorname{cosupp}(\mathfrak a^k,ke+1)
 =\operatorname{cosupp}(\mathfrak a,e+1).
\]
\end{lemma}
\begin{proof}
If every germ of $\mathfrak a_x$ has order at least $e+1$, then every
product of $k$ such germs has order at least $k(e+1)$.  Finite sums and
holomorphic multiples preserve this lower bound.  Since
$ke+1\leq k(e+1)$, every germ of $(\mathfrak a^k)_x$ has order at least
$ke+1$.  Conversely, if a germ $g\in\mathfrak a_x$ has order at most $e$,
then multiplicativity of the order gives
$\operatorname{ord}_x(g^k)\leq ke$.  The germ $g^k$ belongs to
$(\mathfrak a^k)_x$, so this stalk does not have order at least $ke+1$.
These two implications prove the equality.
\end{proof}

\begin{lemma}[Divisibility thresholds for ideal powers]
\label{hir:divisor-power-divisibility-threshold}
\uses{hir:divisor-power-successor-threshold,hir:divisor-exponent}
\lean{Hironaka.HironakaDivisorPowers.isDivNear_idealPow_mul_succ_iff}
\leanok
For a subset $H\subset M$, a point $x\in M$, $k\geq1$ and $e\geq0$,
the ideal $\mathfrak a^k$ has order at least $ke+1$ along $H$ near $x$
if and only if $\mathfrak a$ has order at least $e+1$ along $H$ near $x$.
\end{lemma}
\begin{proof}
Apply Lemma~\ref{hir:divisor-power-successor-threshold} at each point
of $H$ in a neighbourhood of $x$.  The equality of the two cosupports
identifies the two neighbourhood conditions in both directions.
\end{proof}

\begin{lemma}[A power of a nonzero ideal stalk is nonzero]
\label{hir:divisor-power-nonzero}
\uses{hir:order-power,hir:divisor-exponent}
\lean{Hironaka.HironakaDivisorPowers.isNonzeroAt_idealPow}
\leanok
If $\mathfrak a_x$ contains a holomorphic germ which does not vanish
identically near $x$, then $(\mathfrak a^k)_x$ has the same property for
every integer $k\geq0$.
\end{lemma}
\begin{proof}
Choose such a germ $g$.  Its power $g^k$ lies in $(\mathfrak a^k)_x$.
At each point where $g$ is nonzero, $g^k$ is nonzero; these points occur
arbitrarily close to $x$.  This also covers $k=0$, since $g^0=1$.
\end{proof}

\begin{theorem}[Divisor exponents scale with ideal powers]
\label{hir:divisor-exponent-ideal-power}
\uses{hir:divisor-power-divisibility-threshold,hir:divisor-power-nonzero,hir:divisor-exponent}
\dcref{W{\l}odarczyk 2005, Step 2}
\source{wlodarczyk05-hironaka}
\lean{Hironaka.HironakaDivisorPowers.divExp_idealPow}
\leanok
Let $x\in H\subset M$, and let $\mathfrak a$ be an ideal of holomorphic
germs whose stalk at $x$ is nonzero.  For every integer $k\geq0$,
\[
 e_H(\mathfrak a^k,x)=k\,e_H(\mathfrak a,x).
\]
\end{theorem}
\begin{proof}
Put $e=e_H(\mathfrak a,x)$.  Nonvanishing makes this exponent finite:
$\mathfrak a$ has order at least $e$ along $H$ near $x$, and it does not
have order at least $e+1$ there.  Products of germs give divisibility of
$\mathfrak a^k$ to order $ke$, so its exponent is at least $ke$.
For $k\geq1$, Lemma~\ref{hir:divisor-power-divisibility-threshold}
rules out divisibility to order $ke+1$.  Lemma~\ref{hir:divisor-power-nonzero}
ensures that its exponent is also finite, hence is at most $ke$.
For $k=0$, the ideal is the unit ideal and its germ $1$ has order zero
at every point of $H$; its exponent is therefore zero.
\end{proof}

\begin{lemma}[Successor thresholds along a submanifold]
\label{hir:divisor-power-successor-threshold-along}
\uses{hir:divisor-power-successor-threshold,hir:exponent-along,hir:order-basic}
\lean{Hironaka.HironakaDivisorPowers.cosupport_idealPow_mul_succ_eq_along,Hironaka.HironakaDivisorPowers.isDivNearAlong_idealPow_mul_succ_iff}
\leanok
Let $N$ be a submanifold of $M$. For $k\geq1$ and $e\geq0$,
\[
 \operatorname{cosupp}(\mathfrak a^k|_N,ke+1)
 =\operatorname{cosupp}(\mathfrak a|_N,e+1).
\]
Consequently the corresponding neighbourhood divisibility conditions
along $N\cap H$ agree.
\end{lemma}
\begin{proof}
The inclusion from right to left follows by multiplying germs vanishing
to order $e+1$ on $N$, and using $ke+1\leq k(e+1)$.
For the reverse inclusion, fix $x\in N$ and a holomorphic retraction
$r$ onto $N$ near $x$.  A germ $g$ has order at least $q$ on $N$ at
$x$ if and only if $g\circ r$ has ambient order at least $q$.
If $g$ fails to have restricted order $e+1$, then
$g\circ r$ has ambient order at most $e$, and its $k$-th power has
ambient order at most $ke$.  Since
$(g^k)\circ r=(g\circ r)^k$, the germ $g^k\in\mathfrak a^k_x$ fails
to have restricted order $ke+1$.  This proves the reverse inclusion.
Applying the equality pointwise near $x$ proves the divisibility assertion.
\end{proof}

\begin{lemma}[Restricted nonvanishing survives powers]
\label{hir:divisor-power-nonzero-along}
\uses{hir:divisor-power-nonzero,hir:exponent-along}
\lean{Hironaka.HironakaDivisorPowers.isNonzeroAlongAt_idealPow}
\leanok
If $\mathfrak a_x$ contains a holomorphic germ whose restriction does
not vanish identically on $N$ near $x$, then $(\mathfrak a^k)_x$ has
the same property, for every $k\geq0$.
\end{lemma}
\begin{proof}
Choose the germ $g$ in the assumption.  Then $g^k$ belongs to
$(\mathfrak a^k)_x$ and is nonzero at each point of $N$ where $g$ is
nonzero. Such points occur arbitrarily close to $x$.
\end{proof}

\begin{theorem}[Divisor exponents of powers along a submanifold]
\label{hir:divisor-exponent-ideal-power-along}
\uses{hir:divisor-power-successor-threshold-along,hir:divisor-power-nonzero-along,hir:exponent-along}
\dcref{W{\l}odarczyk 2005, Step 2}
\source{wlodarczyk05-hironaka}
\lean{Hironaka.HironakaDivisorPowers.divExpAlong_idealPow}
\leanok
Let $N$ be a submanifold, $x\in N\cap H$, and suppose
$\mathfrak a_x|_N$ is nonzero. For every $k\geq0$,
\[
 e_{H|_N}(\mathfrak a^k,x)=k\,e_{H|_N}(\mathfrak a,x).
\]
\end{theorem}
\begin{proof}
Let $e=e_{H|_N}(\mathfrak a,x)$.  A local holomorphic retraction
onto $N$ and nonvanishing ensure that this exponent is finite.
Products give restricted divisibility of $\mathfrak a^k$ to order $ke$.
For $k\geq1$, Lemma~\ref{hir:divisor-power-successor-threshold-along}
excludes restricted divisibility to order $ke+1$, since it would imply
divisibility of $\mathfrak a$ to order $e+1$. Restricted nonvanishing
survives powers, so the exponent of $\mathfrak a^k$ is finite and equals
$ke$. When $k=0$, its unit germ has order zero on $N$, which gives
the asserted equality.
\end{proof}

\begin{lemma}[An invertible divisor ideal has nonzero stalks]
\label{hir:divisor-ideal-nonzero}
\uses{ana:ideal-sheaves,hir:divisor-exponent}
\lean{Hironaka.HironakaDivisorPowers.isNonzeroAt_of_isInvertible}
\leanok
Every stalk of an invertible divisor ideal contains a holomorphic germ
which does not vanish identically near that point.
\end{lemma}
\begin{proof}
Near a point $x$, choose a principal generator $g$ for the ideal.
Invertibility includes the fact that $g$ vanishes identically on no
nonempty open subset of that neighbourhood. Every open neighbourhood
of $x$ therefore contains a point where $g$ is nonzero. The germ $g$
belongs to the ideal stalk at $x$ and has the asserted property.
\end{proof}

\begin{corollary}[Multiplicities of powers of an actual divisor ideal]
\label{hir:divisor-ideal-power-multiplicity}
\uses{hir:divisor-ideal-nonzero,hir:divisor-exponent-ideal-power}
\lean{Hironaka.HironakaDivisorPowers.divExp_divisor_idealPow}
\leanok
For an effective divisor $D$ on $M$, $x\in H$, and $k\geq0$,
\[
 e_H(\mathcal O(-D)^k,x)=k\,e_H(\mathcal O(-D),x).
\]
\end{corollary}
\begin{proof}
Apply Theorem~\ref{hir:divisor-exponent-ideal-power} to the actual
invertible ideal $\mathcal O(-D)$; its holomorphic nonzero stalk members
are supplied by Lemma~\ref{hir:divisor-ideal-nonzero}.
\end{proof}

\subsection{Removing a genuine connected smooth component}

The next results construct one global multiplicity and a residual
effective divisor from the actual invertible ideal. The component is
assumed connected: a possibly disconnected member of a simple normal
crossings list need not carry one global multiplicity.

\begin{lemma}[Local constancy of an ideal exponent on a smooth carrier]
\label{hir:divisor-exponent-locally-constant-carrier}
\uses{hir:divisor-exponent}
\lean{Hironaka.HironakaDivisorResidual.isLocallyConstant_divExp_on_submanifold}
\leanok
Let $W\subset M$ be a closed submanifold and $\mathfrak a$ a coherent
ideal nonzero at every stalk. The function
$x\mapsto e_W(\mathfrak a,x)$ on $W$ is locally constant.
\end{lemma}
\begin{proof}
At each $x\in W$, the regularity of $W$ and nonvanishing of
$\mathfrak a$ in a neighbourhood give local constancy of the divisor
exponent on $W$. Restrict this neighbourhood to the subspace $W$.
Thus every point of $W$ has a neighbourhood on which the displayed
function equals its value at that point.
\end{proof}

\begin{lemma}[A connected carrier supplies one actual multiplicity]
\label{hir:divisor-exponent-connected-carrier}
\uses{hir:divisor-exponent-locally-constant-carrier,hir:divisor-exponent}
\lean{Hironaka.HironakaDivisorResidual.exists_divExp_on_connected_submanifold}
\leanok
Let $W$ be a nonempty connected closed submanifold of fixed dimension
and $\mathfrak a$ a coherent ideal nonzero at every stalk. There is
$a\in\mathbb N$ such that, for every $x\in W$ and $r\geq0$,
\[
 \mathfrak a\text{ has order at least }r\text{ along }W
 \text{ near }x \quad\Longleftrightarrow\quad r\leq a.
\]
\end{lemma}
\begin{proof}
A locally constant function on a connected space is constant: each
nonempty fibre is both open and closed, so its complement is empty.
Choose $p\in W$ and put $a=e_W(\mathfrak a,p)$. By
Lemma~\ref{hir:divisor-exponent-locally-constant-carrier},
$e_W(\mathfrak a,x)=a$ for every $x\in W$. Nonzero stalks ensure
that the finite exponent characterizes divisibility, proving the
displayed equivalence.
\end{proof}

\begin{lemma}[A smooth hypersurface has empty interior]
\label{hir:divisor-smooth-carrier-empty-interior}
\uses{hir:snc}
\lean{Hironaka.HironakaDivisorResidual.interior_eq_empty_of_smooth_hypersurface}
\leanok
If $W$ is a closed smooth submanifold of dimension $d$ in an
$n$-dimensional complex manifold and $d+1=n$, then $W$ has empty
interior.
\end{lemma}
\begin{proof}
An interior point of $W$ has an open neighbourhood contained in $W$;
in the ambient chart it is therefore regular of dimension $n$.
But every point of $W$ is regular of dimension $d$. The dimension
of a regular germ is unique, so $d=n$, contrary to $d+1=n$.
\end{proof}

\begin{theorem}[Actual residual after removing a connected smooth component]
\label{hir:divisor-connected-component-residual}
\uses{hir:divisor-exponent-connected-carrier,hir:divisor-smooth-carrier-empty-interior,hir:divisor-ideal-nonzero,hir:hadamard-bridge}
\dcref{W{\l}odarczyk 2005, Step 2}
\source{wlodarczyk05-hironaka}
\lean{Hironaka.HironakaDivisorResidual.exists_residual_on_connected_hypersurface}
\leanok
Let $D$ be an effective divisor on an $n$-dimensional complex
manifold $M$, and let $W$ be a nonempty connected closed smooth
submanifold of dimension $d$ with $d+1=n$. There exist $a\geq0$
and an effective divisor $R$ on $M$ such that
\[
 \begin{aligned}
 e_W(\mathcal O(-D),x)&=a &&(x\in W),\\
 \mathcal O(-D)&=\mathcal I_W^a\mathcal O(-R),\\
 \mathcal O(-R)&=(\mathcal O(-D):\mathcal I_W^a),
 \end{aligned}
\]
and $W$ is not contained in the support of $R$.
\end{theorem}
\begin{proof}
The ideal $\mathcal O(-D)$ is coherent and has nonzero stalks.
Lemma~\ref{hir:divisor-exponent-connected-carrier} therefore supplies
the constant finite exponent $a$. At every $x\in W$, all germs of
$\mathcal O(-D)_x$ vanish to order at least $a$ on $W$. The Hadamard
bridge gives the global inclusion
$\mathcal O(-D)\subset\mathcal I_W^a$, including the stalks off $W$.

The smooth hypersurface ideal $\mathcal I_W$ is invertible: locally
on $W$ it is generated by one flattening coordinate, and off $W$
by $1$; its local generator vanishes identically on no open subset
because $W$ has empty interior. Its power $\mathcal I_W^a$ is
invertible as well. The colon ideal
$\mathfrak r=(\mathcal O(-D):\mathcal I_W^a)$ is invertible.
Indeed, for local generators $u$ and $v$ of these two nested
invertible ideals, the inclusion gives $u=qv$ with $q$ holomorphic.
The germ $q$ generates the colon ideal: cancellation by $v$ proves
one inclusion, and multiplication by $v$ proves the other.
Nonvanishing of $u$ on every nonempty open set forces the same
nonvanishing for $q$. Thus $\mathfrak r=\mathcal O(-R)$ for an
effective divisor $R$, and the same cancellation yields
$\mathcal O(-D)=\mathcal I_W^a\mathcal O(-R)$.

If $W$ were contained in the support of $R$, every germ of the
coherent ideal $\mathcal O(-R)$ would vanish on $W$ near its base
point: membership propagates to nearby stalks, and each such germ
vanishes at the points of $W$. Hence
$\mathcal O(-R)\subset\mathcal I_W$, and the product identity would
give $\mathcal O(-D)\subset\mathcal I_W^{a+1}$.
At any $p\in W$ this implies divisibility to order $a+1$, contradicting
the characterization of the exponent $a$. Therefore $W$ is not
contained in the support of $R$.
\end{proof}

\begin{corollary}[Powers of the actual component factorization]
\label{hir:divisor-connected-component-residual-power}
\uses{hir:divisor-connected-component-residual}
\lean{Hironaka.HironakaDivisorResidual.exists_residual_power_on_connected_hypersurface}
\leanok
Under the hypotheses of Theorem~\ref{hir:divisor-connected-component-residual},
for each $k\geq0$ the same actual $a$ and residual $R$ satisfy
\[
 \mathcal O(-D)^k=\mathcal I_W^{ak}\mathcal O(-R)^k.
\]
\end{corollary}
\begin{proof}
Construct $a$ and $R$ by the preceding theorem, then take the $k$-th
power of its whole ideal identity. Products of coherent ideals commute
and associate, so the $k$-th power of a product is the product of the
$k$-th powers, and $(\mathcal I_W^a)^k=\mathcal I_W^{ak}$.
The exponent zero gives the unit ideal on both sides.
\end{proof}

\begin{definition}[Finite product of component powers]
\label{hir:divisor-finite-component-product}
\uses{ana:ideal-sheaves}
\lean{Hironaka.HironakaDivisorResidual.componentPowerProduct}
\leanok
For a finite ordered list $A=((W_1,a_1),\ldots,(W_s,a_s))$ of subsets
of $M$ and nonnegative integers, put
\[
 P(A)=\mathcal I_{W_1}^{a_1}\cdots\mathcal I_{W_s}^{a_s},\qquad
 P(\varnothing)=\mathcal O_M.
\]
\end{definition}

\begin{lemma}[Coherence of a finite component product]
\label{hir:divisor-finite-component-product-coherent}
\uses{hir:divisor-finite-component-product}
\lean{Hironaka.HironakaDivisorResidual.isCoherent_componentPowerProduct}
\leanok
If every $W_i$ is a closed submanifold, then $P(A)$ is coherent.
\end{lemma}
\begin{proof}
The vanishing ideal of each closed submanifold is coherent. Powers
and products of coherent ideals are coherent. Induction on the list,
starting from the coherent unit ideal, proves the claim.
\end{proof}

\begin{theorem}[Extraction along a finite family of connected hypersurfaces]
\label{hir:divisor-finite-connected-component-residual}
\uses{hir:divisor-connected-component-residual,hir:divisor-finite-component-product-coherent}
\lean{Hironaka.HironakaDivisorResidual.exists_residual_on_finite_connected_hypersurfaces}
\leanok
Let $E=(W_1,\ldots,W_s)$ be a finite list of nonempty connected closed
smooth hypersurfaces in $M$, and $D$ an effective divisor. There exist
nonnegative integers $a_1,\ldots,a_s$ and an effective divisor $R$ such that
\[
 \mathcal O(-D)=\mathcal I_{W_1}^{a_1}\cdots
 \mathcal I_{W_s}^{a_s}\mathcal O(-R),
 \qquad \operatorname{Supp}(R)\subset\operatorname{Supp}(D),
\]
and no $W_i$ is contained in $\operatorname{Supp}(R)$.
The list may be empty and may contain repeated members.
\end{theorem}
\begin{proof}
Induct on the length. For the empty list take $R=D$; the product is
the unit ideal. For $E=(W,E')$, Theorem~\ref{hir:divisor-connected-component-residual}
produces $a$ and $D_1$ with
$\mathcal O(-D)=\mathcal I_W^a\mathcal O(-D_1)$ and
$W\not\subset\operatorname{Supp}(D_1)$. The induction hypothesis
applied to $D_1$ and $E'$ produces $A'$ and $R$ with
$\mathcal O(-D_1)=P(A')\mathcal O(-R)$,
$\operatorname{Supp}(R)\subset\operatorname{Supp}(D_1)$, and no
member of $E'$ contained in $\operatorname{Supp}(R)$.
Associativity of coherent ideal products gives
$\mathcal O(-D)=(\mathcal I_W^aP(A'))\mathcal O(-R)$.
A zero of every germ of $\mathcal O(-D_1)$ is a zero of every germ
of its product with $\mathcal I_W^a$, so
$\operatorname{Supp}(D_1)\subset\operatorname{Supp}(D)$.
These inclusions give the asserted support inclusion and also rule
out $W\subset\operatorname{Supp}(R)$.
The members of $E'$ are handled by the induction hypothesis. The
argument does not require the members to be distinct, so repeated
members satisfy the same conclusion.
\end{proof}

\begin{remark}[Scope of the multiplicity calculation]
Each exponent extracted from a single connected component is defined
by neighbourhood divisibility of stalks. In the finite theorem, $a_i$
is the exponent of the intermediate divisor after the preceding members
have been removed; these exponents follow the order of the list. They
are not asserted to be the exponents of the original $D$ for every list
position. In particular a repeated member has exponent zero after its
first removal. For a smooth hypersurface the local exponent agrees with the largest
power of its vanishing ideal dividing the stalk.  Constructing a global
identity $\mathcal O(-D)=\prod_i\mathcal I_{E_i}^{a_i}$ for a general
simple-normal-crossings-supported divisor $D$, and then identifying
$D=\sum_i a_iE_i$, remains an additional obligation.
The finite extraction theorem leaves an actual residual factor. It
does not prove this factor is the unit ideal from an SNC support
assumption, nor does it refine disconnected SNC members into their
connected components.
\end{remark}

\begin{thebibliography}{9}
\bibitem{demailly-cadg} J.-P. Demailly, \emph{Complex analytic and differential geometry}, Chapter II, sections 4--5.
\bibitem{wlodarczyk05-hironaka} J. W{\l}odarczyk, \emph{Simple Hironaka resolution in characteristic zero}, J. Amer. Math. Soc. 18 (2005), 779--822.
\bibitem{bierstone-milman97-invent-math} E. Bierstone and P. D. Milman, \emph{Canonical desingularization in characteristic zero by blowing up the maximum strata of a local invariant}, Invent. Math. 128 (1997), 207--302.
\end{thebibliography}
\end{document}
